\section*{第 \ref{ch:g7:chemical-reactions} 章——化学反应：反应物与生成物}

\begin{solution}{exo:g7:chemical-reactions:1}
\omterm{def:g7:chemical-reactions:physical-transformation}{物理变化}：冰熔化，糖\omterm{def:g2:mixing-and-dissolving:dissolve}{溶解}。\omterm{def:g5:irreversible-changes:chemical-change}{化学变化}：火柴\omterm{def:g4:what-a-fire-needs:burning}{燃烧}，
面包烘烤，铁钉生锈。
\end{solution}

\begin{solution}{exo:g7:chemical-reactions:2}
\omterm{def:g7:chemical-reactions:reactant}{反应物}是被消耗的\omterm{def:g6:pure-substances-and-mixtures:species}{化学物种}；\omterm{def:g7:chemical-reactions:reactant}{生成物}是生成的新
\omterm{def:g6:pure-substances-and-mixtures:species}{化学物种}。
\end{solution}

\begin{solution}{exo:g7:chemical-reactions:3}
碳 $+$ 氧 $\longrightarrow$ 二氧化碳。
\end{solution}

\begin{solution}{exo:g7:chemical-reactions:4}
\omterm{def:g7:chemical-reactions:reactant}{反应物}：铁和氧。\omterm{def:g7:chemical-reactions:reactant}{生成物}：氧化铁。
\end{solution}

\begin{solution}{exo:g7:chemical-reactions:5}
水雾说明有水；澄清石灰水变浑浊，说明有二氧化碳。
\end{solution}

\begin{solution}{exo:g7:chemical-reactions:6}
锌 $+$ 盐酸 $\longrightarrow$ 氯化锌 $+$ 氢气。
\end{solution}

\begin{solution}{exo:g7:chemical-reactions:7}
反应前：4（甲烷\omterm{def:g7:atoms-and-molecules:molecule}{分子}中的）。反应后：$2 + 2 = 4$
（每个水\omterm{def:g7:atoms-and-molecules:molecule}{分子}中各有两个）。
\end{solution}

\begin{solution}{exo:g7:chemical-reactions:8}
没有生成新的\omterm{def:g6:pure-substances-and-mixtures:species}{化学物种}：盐仍然存在，分散在
水中，\omterm{def:g3:separating-mixtures:evaporate}{蒸发}掉水即可重新得到它。这是
\omterm{def:g7:chemical-reactions:physical-transformation}{物理变化}。
\end{solution}

\begin{solution}{exo:g7:chemical-reactions:9}
大部分木材与氧气反应，生成了气体，
即二氧化碳和水蒸气，它们逸入\omterm{def:g4:air-a-mixture-of-gases:air}{空气}中。
只剩下灰烬。
\end{solution}

\begin{solution}{exo:g7:chemical-reactions:10}
氢气 $+$ 氧气 $\longrightarrow$ 水。\omterm{def:g7:chemical-reactions:reactant}{生成物}使
无水硫酸铜变蓝。
\end{solution}

\begin{solution}{exo:g7:chemical-reactions:11}
反应前：$2 \times 2 = 4$ 个氢\omterm{def:g7:atoms-and-molecules:atom}{原子}和 2 个氧\omterm{def:g7:atoms-and-molecules:atom}{原子}。反应后：
两个水\omterm{def:g7:atoms-and-molecules:molecule}{分子}含有 $2 \times 2 = 4$ 个氢\omterm{def:g7:atoms-and-molecules:atom}{原子}和 $2 \times 1 =
2$ 个氧\omterm{def:g7:atoms-and-molecules:atom}{原子}。\omterm{def:g7:atoms-and-molecules:atom}{原子}的种类与数目均不变。
\end{solution}

\begin{solution}{exo:g7:chemical-reactions:12}
氧\omterm{def:g7:atoms-and-molecules:atom}{原子}不能凭空产生：图中反应前有 2 个氧\omterm{def:g7:atoms-and-molecules:atom}{原子}，
反应后有 3 个。正确的写法是：一个碳\omterm{def:g7:atoms-and-molecules:atom}{原子} $+$ 一个氧\omterm{def:g7:atoms-and-molecules:molecule}{分子}
$\longrightarrow$ 一个二氧化碳\omterm{def:g7:atoms-and-molecules:molecule}{分子}，没有剩余。
\end{solution}

\begin{solution}{pb:g7:chemical-reactions:1}
\textbf{1.} 甲烷和氧气。

\textbf{2.} 水。

\textbf{3.} 二氧化碳。

\textbf{4.} 否：它没有被消耗，
也不写在\omterm{def:g7:chemical-reactions:word-equation}{文字表达式}中。

\textbf{5.} 甲烷 $+$ 氧气 $\longrightarrow$ 二氧化碳 $+$
水。

\textbf{6.} \omterm{def:g5:irreversible-changes:chemical-change}{化学变化}：甲烷、氧气两种物种消失，
生成两种新物种（可通过检验鉴别），并放出热量和光。

\textbf{7.} \ce{CH4}, \ce{O2}, \ce{CO2}, \ce{H2O}.

\textbf{8.} 1 个碳\omterm{def:g7:atoms-and-molecules:atom}{原子}、4 个氢\omterm{def:g7:atoms-and-molecules:atom}{原子}、$2 \times 2 = 4$ 个氧\omterm{def:g7:atoms-and-molecules:atom}{原子}。

\textbf{9.} 1 个碳\omterm{def:g7:atoms-and-molecules:atom}{原子}；$2 \times 2 = 4$ 个氢\omterm{def:g7:atoms-and-molecules:atom}{原子}；$2 + 2 \times 1 = 4$
个氧\omterm{def:g7:atoms-and-molecules:atom}{原子}。\omterm{def:g7:atoms-and-molecules:atom}{原子}的种类和数目都相同，
只是重新排列了。

\textbf{10.} $10 \times 2 = 20$ 个氧\omterm{def:g7:atoms-and-molecules:molecule}{分子}。

\textbf{11.} 10 个二氧化碳\omterm{def:g7:atoms-and-molecules:molecule}{分子}。

\textbf{12.} $10 \times 2 = 20$ 个水\omterm{def:g7:atoms-and-molecules:molecule}{分子}。
\end{solution}
